\documentclass[10pt,a4paper]{amsart}
\usepackage[margin=19mm]{geometry}
\usepackage{amsmath,amssymb,mathtools}
\usepackage[scr=rsfs,cal=euler]{mathalfa}
\usepackage{xfrac}
\usepackage[unicode,hidelinks,bookmarks=false]{hyperref}
\newcommand{\ie}{i.e.\ }
\newcommand{\cf}{cf.\ }
\newcommand{\resp}{respectively\ }
\newcommand{\Subsection}{\subsection}
\providecommand{\nobreakdash}{\nobreak\hbox{-}}
\setlength{\emergencystretch}{2em}
\multlinegap0pt
\usepackage{bm}
\usepackage{bbm}
\usepackage{dsfont}
\usepackage{xspace}
\usepackage{cancel}
\usepackage{mathtools}
\usepackage{amsthm}
\newtheorem{lem}{Lem}
\newtheorem{prop}{Prop}
\newtheorem{thm}{Thm}
\newcommand*{\N}{\mathbb{N}}
\newcommand*{\Z}{\mathbb{Z}}
\title[The 1/3-phenomenon of placement probabilities of tilings in ]{The 1/3-phenomenon of placement probabilities of tilings in the semiregular hexagon (Theorem 3)}
\author{Marcus Sch\"onfelder}
\date{Source: arXiv:2607.10449v2 (CC BY 4.0)}
\begin{document}
\maketitle

\noindent\emph{Source and licence.} The 1/3-phenomenon of placement probabilities of tilings in the semiregular hexagon. Version arXiv:2607.10449v2 (CC BY 4.0), distributed under the Creative Commons Attribution 4.0 International licence, as recorded in the arXiv licence metadata for this exact version and shown on the abstract page https://arxiv.org/abs/2607.10449v2. Full rights notice: licenses/arxiv-2607.10449-CC-BY-4.0.txt.\par\smallskip

\noindent\textit{Editorial scope.} Counted unit: the Thm printed as Theorem 3 of the article (source0.tex lines 602--604, source label \texttt{pararedA}), delivered together with the article's own complete original proof of it (source0.tex lines 606--626, one \texttt{proof} environment of 21 lines). No mathematics is written, repaired or re-derived: statement and proof are the article's own text line for line. Required context retained uncounted: 1/3 (lines 228--232); redA (lines 367--373); IFshape (lines 398--406). Every definition, display and label of those lines is retained unchanged. Omitted from this package and not redistributed: 1--227, 233--366, 374--397, 407--601, 605--605, 627--859. Nothing inside a retained excerpt is omitted or altered: every retained body is delivered exactly as the article's lines are written, so no omission receipt of any kind accompanies this package. Citations: the retained excerpts carry no citation command at all, so no bibliography entry of the article is transcribed and no reference is redistributed. Reference adaptation: none was needed and none was made; every reference inside the delivered unit resolves inside it, so no unresolved reference or citation is produced. Printed slips kept literal: no printed word, symbol, hypothesis or number of the counted statement or of its proof was altered. Layout and class: the article's own class is \texttt{amsart} with packages including graphicx, babel, inputenc, tabularx, setspace, subfigure, wrapfig, xcolor, listings, amsmath, amsfonts, amssymb, amsthm, dlfltxbcodetips; the wrapper is the lane's pinned \texttt{amsart} editorial wrapper at 10pt on A4, which restates the theorem-like environments the retained text needs and adds the article's own macro definitions that the retained text needs (\texttt{\textbackslash N}, \texttt{\textbackslash Z}), with extra wrapper packages (bm, bbm, dsfont, xspace, cancel, mathtools). The delivered numbering is the unit's own. Assets: no retained span contains a figure environment, a graphics-inclusion command, a PDF-page inclusion, a table or a TikZ picture; no image file is copied into this package, the package declares no asset dependency and no image file is redistributed with it. Licence: version arXiv:2607.10449v2 of this article is distributed under the Creative Commons Attribution 4.0 International licence, as recorded in the arXiv licence metadata for this exact version and shown on the abstract page https://arxiv.org/abs/2607.10449v2. A full rights notice is delivered at licenses/arxiv-2607.10449-CC-BY-4.0.txt. Characters: every retained excerpt is pure ASCII, so the delivered engine, XeLaTeX with its default Latin Modern text font, reports no missing character.
\begin{thm}[\sc{The $1/3$-phenomenon}]\label{1/3}
    For all $a,b,c,x,y\in\Z$ and all positive integers $n$ with $n\geq N_{a,b,c,x,y}\in\N$ large enough we have
    \[\mathbb{P}(a,b,c,x,y;n)=\frac{1}{3}+f_{a,b,c,x,y}(n)\cdot\binom{2n}{n}^3\bigg/\binom{6n+2}{3n+1}\]
    where $f_{a,b,c,x,y}(n)$ is a rational function in $n$.
\end{thm}

\begin{lem}\label{redA}
    For all $a>0$, all $b,c,x,y$ non-negative integers and $n\in\N$ large enough it holds that
    \begin{align}
        \mathbb{P}(a,b,c,x,y;n)&-\mathbb{P}(a-1,b,c,x,y;n)\notag\\
        &=\frac{(c+2n)!}{(b+2n+1)_{c+2n}}I_{a,b,c,x,y}(a+2n;n)\cdot F_{a,b,c,x,y}(a+2n;n).\notag
    \end{align}
\end{lem}

\begin{prop}\label{IFshape}
    For all fixed $a,b,c,x,y\in \N\cup\{0\}$ and $n\in \N\cup\{0\}$ large enough we have the following.
    \begin{itemize}
        \item[(i)] There exists a rational function $q_{a,b,c,x,y}(n)$ in $n$ such that
        \[I_{a,b,c,x,y}(a+2n;n)=(-1)^n\frac{(3n+1)!}{(n!)^3}\cdot q_{a,b,c,x,y}(n).\]
        \item[(ii)] There exists a rational function $r_{a,b,c,x,y}(n)$ in $n$ such that
        \[F_{a,b,c,x,y}(a+2n;n)=(-1)^n\frac{(2n)!(3n+1)!(4n)!}{(n!)^3 (6n+2)!}\cdot r_{a,b,c,x,y}(n).\]
    \end{itemize}
\end{prop}

\begin{thm}[\sc{Reduction of the parameter $a$}]\label{pararedA}
    We have that $\mathbb{P}(a,b,c,x,y;n)$ with non-negative integer arguments $a,b,c,x,y\geq 0$ satisfies Theorem~\ref{1/3} if $\mathbb{P}(0,b,c,x,y;n)$ fulfils it. 
\end{thm}

\begin{proof}
    By Lemma~\ref{redA} we have
     \begin{align}
        \mathbb{P}(a,b,c,x,y;n)&-\mathbb{P}(a-1,b,c,x,y;n) \label{diffeq}\\
        &=\frac{(c+2n)!}{(b+2n+1)_{c+2n}}I_{a,b,c,x,y}(a+2n;n)\cdot F_{a,b,c,x,y}(a+2n;n).\notag
    \end{align}
    Now, having a closer look at the factors on the right-hand side we see
    \begin{align}
        \frac{(c+2n)!}{(b+1+2n)_{c+2n}}=\frac{(c+2n)!}{\frac{(b+1+c+4n)!}{(b+2n)!}}=\frac{(2n)!}{\frac{(4n)!}{(2n)!}}q_1(n)=\frac{(2n)!(2n)!}{(4n)!}q_1(n)\notag
    \end{align}
    for some rational function $q_1(p)$ since for a positive integer constants $C$ and $k$ we always have $(C+k\cdot n)!=(k\cdot n)!\cdot \text{polynomial}(n)$. Moreover, Proposition~\ref{IFshape} yields that there exist rational functions $q_2(n)$ and $q_3(n)$ such that
    \begin{align}
        &\frac{(c+2n)!}{(b+2n+1)_{c+2n}}I_{a,b,c,x,y}(a+2n;n)\cdot F_{a,b,c,x,y}(a+2n;n)\notag\\
        &=\frac{(2n)!(2n)!}{(4n)!}q_1(n)\cdot (-1)^n\frac{(3n+1)!}{(n!)^3}\cdot q_2(n)\cdot (-1)^n\frac{(2n)!(3n+1)!(4n)!}{(n!)^3(6n+2)!}\cdot q_3(n)\notag\\
        &=\frac{(2n)!(2n)!(2n)!}{(n!)^6}\cdot \frac{(3n+1)!(3n+1)!}{(6n+2)!}\big(q_1(n)\cdot q_2(n)\cdot q_3(n)\big)\notag\\
        &=q(n)\cdot \binom{2n}{n}^3\bigg/\binom{6n+2}{3n+1}\notag
    \end{align}
    where $q(n)=q_1(n)\cdot q_2(n)\cdot q_3(n)$ is again a rational function in $n.$ Hence, if we insert this into Equation~\eqref{diffeq} we instantly obtain
    \[\mathbb{P}(a,b,c,x,y;n)=\mathbb{P}(a-1,b,c,x,y;n)+q(n)\cdot \binom{2n}{n}^3\bigg/\binom{6n+2}{3n+1}\]
    which implies by induction on the parameter $a$ that $\mathbb{P}(a,b,c,x,y;n)$ fulfils Theorem~\ref{1/3} if $\mathbb{P}(0,b,c,x,y;n)$ does.
\end{proof}

\end{document}
